\section{Protocol: Phase-Covariant Transport and Gauge Cost}
\label{sec:protocol}

Without a global phase reference, transport between independently rotated local frames must transform with those frames. \nolinkurl{Protocol} proves that a fixed coefficient invariant under every independent phase rotation vanishes, whereas a transforming link preserves the transport expression. This supplies a conditional reason for a connection; it does not fix the gauge group or its coupling.

\subsection{Links and Minimal Loops}
For an Abelian connection write $U_\mu(x)=e^{-ig\ell A_\mu(x)}$. The observable link phase is unchanged by $g\mapsto g/c$, $A\mapsto cA$ for nonzero $c$ (\nolinkurl{Gauge.phase_normalization_freedom}). Thus a link-speed coefficient is not a determination of $g$.

For a supplied group, a product of links around a closed loop transforms by conjugation at its base point. Lean checks this for four links and shows that a conjugation-invariant score gives a gauge-invariant cost (\nolinkurl{Gauge}). The choice of group, representation, loop set and score remains model data. The Wilson score is a candidate; invariance and positivity do not make it the unique possible class function.

The $D_4$ nearest-neighbour graph has $24$ roots and admits triangles, so it cannot be treated as a square-only graph without extra structure. \nolinkurl{D4Gauge.tri} contains $192$ ordered root pairs with inner product one. Accounting for the six representations of an unoriented triangle across its three vertices gives a triangle density of $32$ per site. The proved count is the ordered-pair count; this normalization must accompany its use in an action.

\subsection{The Actual Cosine Cost and Its Continuum Approximation}
\label{sec:exact_truncation}
For an Abelian loop phase, use $1-\cos\theta$. Within the explicit comparison/repair
contract of the preceding section, this functional form follows from the flag probability:
$W$ refreshes, each requiring $m$ primitive repairs on a flag, cost
$(c_0mW/2)(1-\cos\theta)$. Thus $\beta=c_0mW/2$ within this event convention.
The convention, refresh assignment and repair realization must still be justified for the
physical substrate; symmetry alone does not select this comparison protocol. It is nonnegative and bounded by two (\nolinkurl{Gauge.phase_cost_bounds}). It is not an exact quadratic function at finite phase; \nolinkurl{wilson_not_exactly_quadratic} gives a counterexample. Finite node capacity does not cancel its higher Taylor terms.

The new \nolinkurl{GaugeContinuum} module bounds the approximation. For spacing $h\ne0$ and constant triangle coefficient $a$, define
\begin{equation}
 Q_h(a)=\frac{1-\cos(h^2a)}{h^4}.
\end{equation}
When $|h^2a|\leq1$, Lean proves
\begin{equation}
 |Q_h(a)-a^2/2|\leq\frac5{96}|h|^4|a|^4,
 \qquad \lim_{h\to0,\,h\ne0}Q_h(a)=a^2/2.
\end{equation}
This is a controlled small-spacing result. At fixed nonzero microscopic spacing it describes a long-wavelength approximation, not exact elimination of corrections.

\subsection{The Geometric Quadratic Coefficient on \texorpdfstring{$D_4$}{D4}}
For $B^p_{\mu\nu}=a_\mu b_\nu-a_\nu b_\mu$, the finite triangle sum satisfies
\begin{equation}
 \sum_p B^p_{\mu\nu}B^p_{\rho\sigma}
 =96(\delta_{\mu\rho}\delta_{\nu\sigma}-\delta_{\mu\sigma}\delta_{\nu\rho}).
\end{equation}
Thus for every antisymmetric field $F$,
\begin{equation}
 \sum_p(F\mathbin{\cdot}B^p)^2=192\sum_{\mu,\nu}F_{\mu\nu}^2
\end{equation}
(\nolinkurl{D4Gauge.area_isotropic}, \nolinkurl{flux_sum}). No added squares are needed for this isotropy. The determinant of the Cartan matrix is four; the cell-volume factor used in the coordinate realization is two.

For the supplied constant-field Abelian phase $h^2(g/4)(F\mathbin{\cdot}B^p)$, sixfold triangle accounting and cell volume give
\begin{equation}
 \mathcal S_h(F)=\frac{\beta}{12}\sum_p Q_h\bigl((g/4)(F\mathbin{\cdot}B^p)\bigr).
\end{equation}
Lean proves that the \emph{actual cosine action} tends to
\begin{equation}
 \lim_{h\to0}\mathcal S_h(F)=\frac{\beta g^2}{2}\sum_{\mu,\nu}F_{\mu\nu}^2.
\end{equation}
With $g\ne0$ and $\beta=1/(2g^2)$ this is the canonical Euclidean Maxwell coefficient $1/4$ (\nolinkurl{GaugeContinuum.maxwell_limit}). The normalization identity does not select $\beta$ or a physical gauge coupling.

This completion covers constant Abelian curvature. Nonconstant fields require control of geometric flux errors; non-Abelian fields additionally require path-order/BCH and representation-normalization bounds. A continuation from positive Euclidean cost to a physical Lorentzian theory is a further step, not simply the insertion of a minus sign. Scattering and vertex-response calculations remain open; see \cref{app:impedance_power_transfer}.


\paragraph{The event coefficient in the controlled limit.}
\nolinkurl{CarrierResponse.common_response_limit} substitutes the counted coefficient into
the already proved triangle limit, giving
\begin{equation}
 \lim_{h\to0}\mathcal S_h(F)
 =\frac{c_0mWg^2}{4}\sum_{\mu,\nu}F_{\mu\nu}^2.
\end{equation}
This ties the coefficient to a primitive toll, repair multiplicity and refresh count in one
specified contract. Canonical coefficient $1/4$ requires $c_0mWg^2=1$ in these conventions.
For matrix-valued holonomies, representation traces and indices must also be included;
this Abelian result does not determine the physical relative gauge normalization or
which resource each gauge sector refreshes.
